\input{LectureSlides/includes/metropolis_preamble}

% Title page info
\title[Nash Equilibrium]{Chapter 2 (Part 2): Auctions, Iterated Dominance, \& Nash Equilibrium}
\author[ECON 5001]{ECON 5001\\Prof. Paul J. Healy\\Based on: \textit{Game Theory} by Giacomo Bonanno} \color{metrop}
\institute[OSU]{}
\date[]{\vfill {\tiny Updated \today\ at\ \DTMcurrenttime}}

\begin{document}

\frame{\maketitle}

\begin{frame}{Where We Left Off}
We can now:
\begin{itemize}
    \item separate a \textbf{frame} from a \textbf{game}
    \item write preferences as a complete, transitive $\succsim_i$
    \item represent $\succsim_i$ with an ordinal utility function
    \item compare two strategies of one player: \textbf{dominance}
\end{itemize}
\bsk
And we found: if you have a dominant strategy, play it.
\vfill
In this lecture: what if you don't?
\vspace{-0.5in}
\end{frame}

% ==============================================================
% MOBLAB #1 --- PLAY BEFORE ANY THEORY
% ==============================================================

% ---- MobLab run: sealed-bid auction ---------------------------------
% Game: "Private Value Sealed" (Auction topic). Run TWO rounds:
%   Round 1: winner pays own bid   (first price)
%   Round 2: winner pays 2nd-highest bid (second price)
% Same values both rounds if possible, so the bid-shading comparison is clean.
% Do NOT explain the pricing rules beyond what's on the screen.
% ---------------------------------------------------------------------
\begin{frame}{Log In: Auction Game}
\begin{itemize}
    \item MobLab: Two auctions. You're bidders.
    \item You'll be told what the item is worth \textbf{to you}
    \begin{itemize}
        \item That number is yours alone; others have different ones
        \item Random from \$10 to \$100 (equally likely)
    \end{itemize}
    \item Groups of SIX bidders.
    \item Submit one sealed bid. Highest bid wins.
    \begin{itemize}
        \item First game: Winner pays their bid
        \item Second game: Winner pays highest \textit{losing} bid (2nd highest bid)
    \end{itemize}
    \item We'll play it \textbf{twice}, under two different rules
\end{itemize}
\bsk
Your payoff is your value minus what you pay, or zero if you don't win.
\vfill
\centering
No talking. Bid however you think is smart.
\vspace{-0.3in}
\end{frame}

\begin{frame}{What Did You Just Play?}
\textbf{Round 1, first-price auction:} the winner pays \textit{their own bid}.
\bsk
\textbf{Round 2, second-price auction:} the winner pays \textit{the 2nd highest bid}.
\bsk
\vfill
In which round did you bid closer to your true value?\\
Did you (bid $<$ value) or (bid $>$ value)?
\vspace{1.1in}
\end{frame}

% ==============================================================
% 2.3 SECOND-PRICE AUCTION
% ==============================================================

\begin{frame}[standout]
Second-Price Auctions
\end{frame}

\begin{frame}{The Rules}
A \textbf{second-price} (or \textbf{Vickrey}) auction:
\begin{itemize}
    \item sealed bids: nobody sees anyone else's bid
    \item highest bidder wins
    \item the winner pays the \textbf{second-highest} bid, not their own
    \item ties are broken by a rule fixed in advance
\end{itemize}
\bsk
Why would a seller ever do this? It looks like giving money away.
\vfill
\vspace{0.9in}
\end{frame}

\begin{frame}{An Example: The Game-Day Permit}
Suppose OSU auctions off a food truck permit near the stadium for game day.\\
Dos Hermanos and Pitabilities both want it.
\bsk
\begin{itemize}
    \item Dos Hermanos (Player 1): estimates \textbf{\$30k} extra profit if they get it\\
    This is their \textbf{value} for the permit. Write $v_1=30$ (in thousands)
    \item Suppose bids must be one of \$10k, \$20k, \$30k, \$40k, \$50k, \$60k\\
    So $S_1=S_2=\{10,20,30,40,50,60\}$. They have 5 possible strategies.\\
    Write bids as $b_1\in S_1$ and $b_2\in S_2$
    \item \textbf{Second-price auction:}highest bid wins, pays \textit{2nd highest bid}
    \item Ties go to Pitabilities (decided by a coin flip beforehand)
\end{itemize}
\bsk
Dos Hermanos's payoff in the Second Price Auction:
$$\pi_1 = \begin{cases} v_1 - b_2 & \text{if}\ b_1>b_2 \\ 0 & \text{if}\ b_1\leq b_2\end{cases}$$
Assume utility $=$ money (selfish \& greedy)
\end{frame}

\begin{frame}{Only Dos Hermanos's Payoffs Shown}
\centering
{\small
\begin{matrixgame}[top={10, 20, 30, 40, 50, 60}, left={10, 20, 30, 40, 50, 60}, player names={Dos Hermanos ($v_1=30$), Pitabilities}]
  \payoffmatrix{0,\cdot & 0,\cdot & 0,\cdot & 0,\cdot & 0,\cdot & 0,\cdot \\
                20,\cdot & 0,\cdot & 0,\cdot & 0,\cdot & 0,\cdot & 0,\cdot  \\
                20,\cdot & 10,\cdot & 0,\cdot & 0,\cdot & 0,\cdot & 0,\cdot  \\
                20,\cdot & 10,\cdot & 0,\cdot & 0,\cdot & 0,\cdot & 0,\cdot  \\
                20,\cdot & 10,\cdot & 0,\cdot & -10,\cdot & 0,\cdot & 0,\cdot \\
                20,\cdot & 10,\cdot & 0,\cdot & -10,\cdot & -20,\cdot & 0,\cdot }
  \foreach \i in {1,...,6}{ \shaderow[red!30]{\i}{\i}{6} }
  \foreach \i in {2,...,6}{ \shaderow[green!30]{\i}{1}{\i-1} }
\end{matrixgame}}
\bsk
\raggedright
Recall: $\pi_1 = \begin{cases} {\color{green!60!black} v_1 - b_2} & \text{if}\ {\color{green!60!black}b_1>b_2} \\ \color{red}0 & \text{if}\ \color{red}b_1\leq b_2\end{cases}$\hspace{0.25in} with $v_1=30$
\end{frame}

\begin{frame}{Reading the Table}
\begin{itemize}
    \item $b_1=30$ is a \underline{\hspace{2in}} strategy \vspace{0.3in}
    \item $b_1=40$ is a \underline{\hspace{2in}} strategy \vspace{0.3in}
    \item Strategies $b_1=30$ and $b_1=40$ are \underline{\hspace{2in}} \vspace{0.3in}
    
    \item Overbidding ($b_1=50$) is weakly dominated by $b_1=30$. Where is it strictly worse? \vspace{0.4in}
    \item Underbidding ($b_1=20$) is weakly dominated by $b_1=30$. Where is it strictly worse?\vspace{0.3in}
\end{itemize}
\end{frame}

\begin{frame}{Vickrey's Theorem}
\textbf{Theorem (Vickrey, 1961).} In a second-price auction, if player $i$ is selfish and
greedy, then bidding their true value ($b_i = v_i$) is a \textbf{weakly dominant strategy}.
\vfill
Remarkable, because:
\begin{itemize}
    \item it holds no matter how many bidders there are
    \item it holds no matter what anyone else bids
    \item bidders don't need to guess, predict, or model their rivals at all!
\end{itemize}
\vfill
\textbf{Lab result:} many people don't realize this, end up overbidding ($b_i > v_i$)\\
``Failure of contingent thinking''
\vspace{-0.4in}
\end{frame}

\begin{frame}{Why It Works: Two Cases}
Let $v_1$ be my value and $b_2$ the other bid. Compare bidding $v_1$ to anything else.
\bsk
\textbf{Case 1: $b_2\le v_1$.} \ Bidding $v_1$ wins and pays $b_2$, for a payoff of $v_1-b_2\ge 0$.
\begin{itemize}
    \item Bidding higher: still win, still pay $b_2$. \underline{\hspace{1.5in}}
    \item Bidding below $b_2$: lose, payoff 0. \underline{\hspace{1.5in}}
\end{itemize}
\bsk
\textbf{Case 2: $b_2>v_1$.} \ Bidding $v_1$ loses, for a payoff of 0.
\begin{itemize}
    \item Bidding anything below $b_2$: still lose. \underline{\hspace{1.5in}}
    \item Bidding at or above $b_2$: win, pay $b_2$. \underline{\hspace{1.5in}}
\end{itemize}
\bsk
Raising your bid can only ever change \textit{whether} you win, never \textit{what you pay}.
\end{frame}

\begin{frame}{The Fine Print}
The theorem assumes the bidder is \textbf{selfish and greedy}.
\bsk
Drop that and it fails. Suppose I also care what \textit{you} pay:
\begin{itemize}
    \item \textbf{Spiteful}: if I lose, I'd rather you paid more \vspace{0.35in}
    \item \textbf{Benevolent}: if I lose, I'd rather you paid less \vspace{0.35in}
\end{itemize}
\bsk
In either case, bidding my value is no longer dominant.
\vfill
Again: \textit{preferences matter!}
\vspace{-0.3in}
\end{frame}

% \begin{frame}{Back to Your Bids}
% Under second-price, truthful bidding is weakly dominant.
% \bsk
% Under first-price, it is not\ldots\ winning means paying your own bid, so bidding your
% value guarantees a payoff of exactly zero. You must \textbf{shade} down.
% \vfill
% \centering
% \begin{tabular}{lc}
% Round 1 (first-price), average bid/value: & \underline{\hspace{1in}} \\[0.3in]
% Round 2 (second-price), average bid/value: & \underline{\hspace{1in}} \\[0.3in]
% \end{tabular}
% \vfill
% \end{frame}

% ==============================================================
% 2.4 PIVOTAL MECHANISM
% ==============================================================

\begin{frame}[standout]
The Pivotal Mechanism
\end{frame}

\begin{frame}{The Pivotal Mechanism}
I'm skipping this due to time.\\
This will \textbf{not} appear on a quiz or exam.\\
In a nutshell:
\begin{itemize}
    \item We're thinking about building a park. 
    \item Person $i$'s value: $v_i$.\hspace{0.5in}Total cost: $C$
    \item Social Welfare: $W = \sum_i v_i - C$
    \item Pareto optimal: Build if $W\geq 0$
    \item Suppose your tax is based on $v_i$. You'd lie about $v_i$!!
\end{itemize}
The fix: the \textbf{pivotal} mechanism
\begin{itemize}
    \item Social Welfare if $i$ didn't exist: $W_{-i}$
    \item Player $i$ is \textbf{pivotal} if $W_{-i}<0$ but $W\geq 0$ (or vice-versa)
    \item Player $i$'s tax is $W_{-i}$
    \item \textbf{Result:} It's a weakly dom. strategy to report $v_i$ truthfully
\end{itemize}
Exactly like a 2nd-price auction!!
\end{frame}

% \begin{frame}{A Real Headline}
% \textit{Davis Enterprise}, January 12, 2001:
% \bsk
% \begin{quote}
% ``By consensus, the Davis City Council agreed Wednesday to order a community-wide
% public opinion poll to gauge how much Davis residents would be willing to pay for a
% park tax and a public safety tax.''
% \end{quote}
% \bsk
% Will people answer honestly?
% \begin{itemize}
%     \item If my answer doesn't affect my tax bill: \vspace{0.35in}
%     \item If my answer does affect my tax bill: \vspace{0.35in}
% \end{itemize}
% \end{frame}

% \begin{frame}{The Setup}
% A public project such as building the Upper Arlington Community Center costs
% \$$C$ in total. There are $n$ residents.
% \bsk
% \begin{itemize}
%     \item Resident $i$'s share of the cost is $c_i$, with $c_1+\cdots+c_n=C$ \\
%       {\footnotesize (shares are announced in advance and need not be equal)}
%     \item Resident $i$ starts with wealth $m_i$
%     \item If built, resident $i$ gets benefit worth $v_i$ (possibly \textbf{negative}) \\
%       {\footnotesize (you live on the street it's going on: traffic, lights, parking)}
% \end{itemize}
% \bsk
% $$U_i(\$m)=\begin{cases} m & \text{if the project is not built}\\ m+v_i & \text{if it is}\end{cases}$$
% \vfill
% \textbf{Efficiency:} build it exactly when the project is worth more to the residents
% together than it costs to build. Social benefit beats total cost.
% \bsk
% $$\underbrace{\textstyle\sum_{i=1}^n v_i}_{\text{social benefit}}\ \ >\ \ \underbrace{C\vphantom{\textstyle\sum}}_{\text{total cost}}$$
% \vspace{-0.4in}
% \end{frame}

% \begin{frame}{Why Efficiency Isn't Enough}
% $n=2$, \ $m_1=50$, \ $m_2=60$, \ $v_1=19$, \ $v_2=-15$, \ $C=6$, \ $c_1=6$, \ $c_2=0$.
% \bsk
% $\sum v_i = 19-15 = 4 < 6 = C$, so it should \textit{not} be built.
% \bsk
% \centering
% \begin{tabular}{l|c|c|}
% \cline{2-3}
%  & Not built & Built \\
% \hline
% \multicolumn{1}{|l|}{Resident 1} & $50$ & $50+19-6=63$ \\
% \hline
% \multicolumn{1}{|l|}{Resident 2} & $60$ & $60-15=45$ \\
% \hline
% \end{tabular}
% \bsk
% \raggedright
% Gain of 13 to one, loss of 15 to the other.
% \vfill
% The government would happily follow the rule if it knew $v_1$ and $v_2$.
% \bsk
% It doesn't. And asking is exactly the problem.
% \vspace{-0.3in}
% \end{frame}

% \begin{frame}{Clarke's Idea}
% Each resident $i$ reports a number $w_i$, which represents their reported benefit (or harm, if negative).\\
% They may lie: nothing forces $w_i=v_i$.
% \bsk
% \textbf{Decision rule:}
% $$\text{Build?}\quad\begin{cases} \text{Yes} & \text{if } \sum_{j=1}^n w_j > C\\[2pt]
% \text{No} & \text{if } \sum_{j=1}^n w_j \le C\end{cases}$$
% \bsk
% \textbf{Then the twist.} Resident $i$ is \textbf{pivotal} if the decision \textit{without their report}
% would differ from the decision \textit{with} it.
% \begin{itemize}
%     \item Not pivotal $\Rightarrow$ they pay no tax
%     \item Pivotal $\Rightarrow$ they pay $\Big|\sum_{j\ne i}w_j-\sum_{j\ne i}c_j\Big|$
% \end{itemize}
% \end{frame}

% \begin{frame}{Worked Example}
% $n=3$, \ $C=10$, \ $c_1=3$, $c_2=2$, $c_3=5$. \ Reports: $w_1=-1$, $w_2=8$, $w_3=3$.
% \bsk
% $\sum w_j=10$, which is \textit{not} greater than $C=10$, so the project is \textbf{not built}.
% \bsk
% \centering
% {\small
% \begin{tabular}{|c|c|c|c|c|}
% \hline
% $i$ & $\sum_{j\ne i}w_j$ & $\sum_{j\ne i}c_j$ & Decision without $i$ & Tax \\
% \hline
% $1$ & $8+3=11$ & $2+5=7$ & Build (pivotal) & $|11-7|=4$ \\
% \hline
% $2$ & $-1+3=2$ & $3+5=8$ & Don't (not pivotal) & $0$ \\
% \hline
% $3$ & $-1+8=7$ & $3+2=5$ & Build (pivotal) & $|7-5|=2$ \\
% \hline
% \end{tabular}}
% \bsk
% \raggedright
% \vfill
% Note residents 1 and 3 each flipped the outcome, and each pays for it.
% \vspace{-0.3in}
% \end{frame}

% \begin{frame}{``Then I'll Just Avoid Being Pivotal''}
% You can guarantee it: report $w_i=c_i$ and you are never pivotal, never taxed.
% \bsk
% Is that smart? Take $n=4$, $C=15$, $c=(5,0,5,5)$. You are resident 1, with
% $m_1=40$ and $v_1=25$. You expect $w_2=-40$, $w_3=15$, $w_4=20$.
% \bsk
% \textbf{Report $w_1=c_1=5$:} \ $\sum w_j = 0 < 15$, not built, and you pay nothing.
% $$U_1 = 40$$
% \bsk
% \textbf{Report truthfully, $w_1=25$:} \ $\sum w_j = 20 > 15$, so it \textit{is} built.
% You are pivotal, taxed $|(-5)-(10)|=15$.
% $$U_1 = 40 + 25 - 5 - 15 = 45$$
% \vspace{-0.3in}
% \end{frame}

% \begin{frame}{Clarke's Theorem}
% \vfill
% \centering
% \textbf{Theorem (Clarke, 1971).} In the pivotal mechanism, reporting truthfully
% ($w_i=v_i$) is a \textbf{weakly dominant strategy} for every player.
% \vfill
% \raggedright
% Same shape as Vickrey: the tax you pay is built out of \textit{other people's} numbers,
% so your report moves the decision without moving your own bill.
% \vfill
% \vspace{0.5in}
% \end{frame}

% ---- Historical figure slide: three portraits in chronological order ----
\begin{frame}{Intellectual History}
\centering
\begin{tabular}{c@{\hspace{0.45in}}c@{\hspace{0.45in}}c}
\includegraphics[height=1.05in]{LectureSlides/graphics/history/VickreyMug.png} &

\includegraphics[height=1.05in]{LectureSlides/graphics/history/ClarkeMug.png} &

\includegraphics[height=1.05in]{LectureSlides/graphics/history/GrovesMug.png} \\[0.04in]
{\small Vickrey, 1961} & {\small Clarke, 1971} & {\small Groves, 1973}
\end{tabular}
\bsk
\raggedright
\begin{itemize}
    \item \textbf{1973}: Groves publishes the pivotal mechanism
    \item \textbf{1971}: Clarke had already published a special case
    \item \textbf{1961}: People realized it's just like Vickrey's 2nd price auction
\end{itemize}
\bsk
The ``pivotal mechanism'' is often called ``the \textbf{VCG} mechanism.''
\vfill
{\footnotesize Vickrey won the Nobel Prize in 1996 but died three days after it was announced.}
\vspace{-0.2in}
\end{frame}

% ==============================================================
% 2.5 ITERATED DELETION
% ==============================================================

\begin{frame}[standout]
Next Solution Concept:\\Iterated Deletion of Dominated Strategies
\end{frame}

\begin{frame}{An Example}
\centering
\begin{matrixgame}[top={L, C, R}, left={T, M, B}, player names={Player 1, Player 2}, br=<2->]
  \payoffmatrix{4,3 & 5,1 & 6,2 \\ 2,1 & 8,4 & 3,6 \\ 3,0 & 9,6 & 2,8}
  \strike<3->{v}{2}
  \strike<4->{h}{2}
  \strike<4->{h}{3}
  \strike<5->{v}{3}
\end{matrixgame}
\bsk
\raggedright
\begin{enumerate}
    \item Underline best-response payoffs
    \item<2-> There are no dominant strategies. But are any \textit{strictly dominated}?
    \begin{itemize}
        \item<2-> Fact: dominated $\Rightarrow$ never a best response (not underlined)
        \item<2-> So, check $M$ and $C$
    \end{itemize}
    \item<2-> $C$ is strictly dominated by \underline{\hspace{0.5in}}. Delete it!!
    \begin{itemize}
        \item Player 1 knows Player 2 will \textbf{not} play $C$
    \end{itemize}
    \item<3-> But now $M$ and $B$ are strictly dominated by \underline{\hspace{0.5in}}. Delete them
    \item<4-> Now $R$ is dominated by $L$. Delete $R$!
    \item<5-> Our solution: $s^*=(T,L)$
\end{enumerate}
\end{frame}

\begin{frame}{IDSDS}
    Our new solution concept:
    \vfill
    {\Large
    \textbf{I}terated\\
    \textbf{D}eletion of\\
    \textbf{S}trictly\\
    \textbf{D}ominated\\
    \textbf{S}trategies}\\
    \vfill
    (Also known as I\textbf{E}SDS, for ``Elimination'' instead of ``Deletion'')
\end{frame}

\begin{frame}{IDSDS}
\textbf{Iterated Deletion of Strictly Dominated Strategies.}
\bsk
Start with game $G$.
\begin{itemize}
    \item $G^1$: delete every strictly dominated strategy, for every player
    \item $G^2$: do it again, inside $G^1$
    \item \ \ $\vdots$
    \item $G^\infty$: the output, reached in finitely many steps if $G$ is finite
\end{itemize}
\bsk
If $G^\infty$ is a single strategy profile, we call it the
\textbf{iterated strict dominant-strategy profile}.\\
Otherwise it's just ``the output of IDSDS.''
\end{frame}

\begin{frame}{A Rationality Logic}
We can justify IDSDS via ``common knowledge of rationality''
\begin{description}
    \item[1.] I know you're rational $\Rightarrow$ you won't play a strictly dominated strategy. $G^1$
    \item[2.] I know you know I'm rational $\Rightarrow$ I know you'll eliminate my remaining strictly dominated strategies. $G^2$
    \item[3.] I know you know I know you're rational $\Rightarrow$ I know you know that I'll eliminate your remaining strictly dominated strategies. $G^3$
    \item[] $\cdots$
    \item[$\infty$.] I know that you know that I know that... $\Rightarrow$ IDSDS. $G^\infty$
\end{description}
    
\end{frame}

\begin{frame}{Shortcoming of IDSDS}
\begin{center}
\begin{matrixgame}[top={W, X, Y, Z}, left={A, B, C, D}, player names={Player 1, Player 2},
                   w=1.4, h=0.65, br=<2->]
  \payoffmatrix{3,4 & 6,6 & 2,1 & 2,0 \\
                2,1 & 1,2 & 5,5 & 2,3 \\
                7,3 & 3,4 & 1,3 & 2,2 \\
                2,8 & 1,1 & 1,2 & 9,0}
  \strike<3->{v}{4}
  \strike<4->{h}{4}
  \strike<5->{v}{1}
  \strike<6->{h}{3}
\end{matrixgame}
\end{center}
\small
\begin{enumerate}
    \item Underline best-response payoffs
    \item<2-> No dominant strategies. But are any \textit{strictly dominated}?
    \begin{itemize}
        \item<2-> Dominated $\Rightarrow$ never a best response. Only $Z$ has no underline
    \end{itemize}
    \item<2-> $Z$ is strictly dominated by \underline{\hspace{0.4in}}. Delete it!
    \item<3-> $D$'s only underline was in $Z$. Now it is dominated by \underline{\hspace{0.4in}}. Delete it!
    \item<4-> $W$'s only underline was in $D$. Now it is dominated by \underline{\hspace{0.4in}}. Delete it!
    \item<5-> $C$'s only underline was in $W$. Now it is dominated by \underline{\hspace{0.4in}}. Delete it!
    \item<6-> Nothing left to delete: $G^\infty$ is a $2\times 2$ game. IDSDS is not \textbf{decisive}.
\end{enumerate}
\vspace{-0.3in}
\end{frame}

% \begin{frame}{Bonanno's Example}
% Four rounds instead of three.
% \bsk
% \centering
% \begin{matrixgame}[top={e, f, g, h}, left={A, B, C, D}, player names={Player 1, Player 2}]
%   \payoffmatrix{6,3 & 4,4 & 4,1 & 3,0 \\ 5,4 & 6,3 & 0,2 & 5,1 \\ 5,0 & 3,2 & 6,1 & 4,0 \\ 2,0 & 2,3 & 3,3 & 6,1}
% \end{matrixgame}
% \bsk
% \raggedright
% \begin{enumerate}
%     \item delete $h$ \ (dominated by $g$) \hspace{0.4in} 3. delete $g$ \ (dominated by $f$)
%     \item delete $D$ \ (dominated by $C$) \hspace{0.35in} 4. delete $C$ \ (dominated by $A$)
% \end{enumerate}
% \bsk
% Output: $\{(A,e),(A,f),(B,e),(B,f)\}$, not a single profile.
% \vspace{-0.3in}
% \end{frame}

% ==============================================================
% MOBLAB #2: BEAUTY CONTEST, PLAYED AFTER IDSDS
% ==============================================================

% ---- MobLab run: (Keynesian) Beauty Contest -------------------------
% Set P = 2/3, integers 1-100. Run 2-3 rounds, showing the winning number
% after each. Deliberately placed AFTER the IDSDS section: students have
% the vocabulary by now, so the payoff is that they still do not play 1 in
% round 1. That gap between knowing the logic and using it is the point,
% and it is what the 'Back to Your Numbers' frame picks up next.
% ---------------------------------------------------------------------
\begin{frame}{Log In Again: Pick a Number}
MobLab: ``Beauty contest''
\vfill
\centering
\large
Choose an integer from \textbf{1} to \textbf{100}.
\bsk
\normalsize
Whoever is closest to \ $\tfrac{2}{3}\times$(the average of everyone's number) \ wins.
\vfill
\raggedright
We'll play a few rounds.
\bsk
\vspace{0.8in}
\end{frame}

% \begin{frame}{Results}
% \centering
% \begin{tabular}{lcc}
%  & \textbf{Average} & \textbf{Winning number} \\[0.1in]
% Round 1: & \underline{\hspace{0.8in}} & \underline{\hspace{0.8in}} \\[0.3in]
% Round 2: & \underline{\hspace{0.8in}} & \underline{\hspace{0.8in}} \\[0.3in]
% Round 3: & \underline{\hspace{0.8in}} & \underline{\hspace{0.8in}} \\[0.3in]
% \end{tabular}
% \bsk
% \raggedright
% \vfill
% What happened between rounds? And why?
% \vspace{-0.3in}
% \end{frame}

\begin{frame}{Back to Your Numbers}
$\tfrac{2}{3}$ of the average can never exceed $\tfrac{2}{3}\times 100\approx 66.7$.
\begin{itemize}
    \item So $67$ strictly dominates everything above it. Delete $68$--$100$. \vspace{0.25in}
    \item Now the target can't exceed $\tfrac{2}{3}\times 67\approx 44.7$. Delete $45$--$67$. \vspace{0.25in}
    \item Now it can't exceed $\tfrac{2}{3}\times 45=30$. Delete $31$--$45$. \vspace{0.25in}
    \item Keep going \ldots \ where does it end? \vspace{0.25in}
\end{itemize}
\vfill
\textbf{Help me:} Why does IDSDS fail? Logically it makes so much sense...
\vspace{-0.3in}
\end{frame}

\begin{frame}{IDWDS, and Why It's Trickier}
What if we delete \textbf{weakly} dominated strategies too?
\bsk
\centering
\begin{matrixgame}[top={L, R}, left={A, T, M, B}, player names={Player 1, Player 2}]
  \payoffmatrix{4,0 & 0,0 \\ 3,2 & 2,2 \\ 1,1 & 0,0 \\ 0,0 & 1,1}
\end{matrixgame}
\bsk
\raggedright
Two perfectly reasonable places to start:
\begin{itemize}
    \item $M$ is strictly dominated by $T$ \ \ \ \ \ \only<1>{$\leftarrow$ delete this first}
    \item $B$ is strictly dominated by $T$ \ \ \ \ \ \only<2>{$\leftarrow$ delete this first}
\end{itemize}
\end{frame}

\begin{frame}{IDWDS}
    Does IDWDS have the same rationality foundation?
    \bsk
    \begin{center}
        \begin{matrixgame}[top={C,D},left={C,D},player names={Sarah,Steven}]
            \payoffmatrix{50,50 & 0,100\\100,0 & 0,0}
        \end{matrixgame}
    \end{center}
    \bsk
    $D$ is weakly dominated but is it \textit{definitely} irrational?
    \bsk
    What if Sarah were \textit{positive} that Steven would play $D$? 
\end{frame}

% ==============================================================
% 2.6 NASH EQUILIBRIUM
% ==============================================================

\begin{frame}[standout]
Our Main Solution Concept:\\
Nash Equilibrium
\end{frame}

\begin{frame}{Still Not Enough}
Games solved outright by IDSDS or IDWDS are the exception.
\bsk
We need something that applies to \textit{every} game.
\vfill
The idea: stop asking ``what will a rational player do?'' and start asking
\bsk
\centering
\large
``which profiles are \textbf{stable}?''
\vfill
\end{frame}

\begin{frame}{An Example}
\begin{center}
    \begin{matrixgame}[top={L,R}, left={U,D}, player names={P1,P2}]
        \payoffmatrix{ 1,1 & 0,0\\ 0,0 & 2,2}
    \end{matrixgame}
\end{center}
Profile $(U,L)$ is \textbf{stable}:
\begin{enumerate}
    \item Player 1 is hurt by deviating from $U$ to $D$
    \item Player 2 is hurt by deviating from $L$ to $R$
\end{enumerate}
\bsk
\onslide<2->{Profile $(D,R)$ is also \textbf{stable}:
\begin{enumerate}
    \item Player 1 is hurt by deviating from $D$ to $U$
    \item Player 2 is hurt by deviating from $R$ to $L$
\end{enumerate}}
    
\end{frame}

\begin{frame}{Definition: Nash Equilibrium (Two Players)}
In a 2-player game, a strategy profile $s^*=(s_1^*,s_2^*)$ is a \textbf{Nash equilibrium} if
\bsk
\begin{enumerate}
    \item for every $s_1\in S_1$: \hspace{0.3in} $\pi_1(s_1^*,s_2^*)\ \ge\ \pi_1(s_1,s_2^*)$
    \bsk
    \item for every $s_2\in S_2$: \hspace{0.3in} $\pi_2(s_1^*,s_2^*)\ \ge\ \pi_2(s_1^*,s_2)$
\end{enumerate}
\vfill
In words: \textbf{assuming the other player sticks to the Nash equilibrium}, neither player can do better by changing their strategy.
\bsk
Does not allow \textbf{joint deviations} (see previous example)
\end{frame}

% ---- Historical figure slide: portrait left, facts right ------------
\begin{frame}{John Forbes Nash Jr. (1928--2015)}
\begin{columns}
\begin{column}{0.32\textwidth}
    \begin{center}
        \includegraphics[height=1.3in]{LectureSlides/graphics/history/NashMug.png}
    \end{center}
\end{column}
\begin{column}{0.68\textwidth}
\begin{itemize}
    \item Duffin's reference letter, 1948, in full: \\
      ``He is a mathematical genius.''
    \item Princeton PhD, 1950: 26 pages, two citations
    \item One of the two was his own
    \item Nobel, 1994, with Harsanyi and Selten
    \item Abel Prize, 2015: PDEs, not games
    \item Schizophrenia from 1959, three decades away
    \item Died May 2015, coming home from the Abel ceremony
\end{itemize}
\end{column}
\end{columns}
\vfill
\vspace{0.5in}
\end{frame}

\begin{frame}{Another Example}
\centering
\begin{matrixgame}[top={L, C, R}, left={T, M, B}, player names={Player 1, Player 2}]
  \payoffmatrix{3,2 & 0,0 & 1,1 \\ 3,0 & 1,5 & 4,4 \\ 1,0 & 2,3 & 3,0}
\end{matrixgame}
\bsk
\raggedright
Verify $(T,L)$ is a Nash equilibrium (NE):
\begin{itemize}
    \item Player 1, holding $L$ fixed: \ $T$ gives 3; \ $M$ gives 3; \ $B$ gives 1 \ \checkmark
    \begin{itemize}
        \item Nothing is \textbf{strictly} better. (Ties are okay.)
    \end{itemize}
    \item Player 2, holding $T$ fixed: \ $L$ gives 2; \ $C$ gives 0; \ $R$ gives 1 \ \checkmark
\end{itemize}
\bsk
So $(T,L)$ is a Nash equilibrium. Find the other one. \vspace{0.5in}
\end{frame}

\begin{frame}{Four Ways to Say It}
\textbf{No regret.} After seeing what the other player did, nobody wishes they had
chosen differently.
\bsk
\textbf{Self-enforcing agreement.} If the players could talk beforehand and shake on
$s^*$, nobody would want to break it, assuming the other keeps it.
\bsk
\textbf{Viable recommendation.} If a third party publicly recommended $s^*$, nobody
would want to deviate, assuming the others comply.
\bsk
\textbf{Rationality with correct beliefs.} Each player best-responds to what they
believe the other will do, and those beliefs are right.
\vfill
\end{frame}

\begin{frame}{Best Replies}
Strategy $s_i$ is a \textbf{best reply} to $s_{-i}$ if
$$\pi_i(s_i,s_{-i})\ \ge\ \pi_i(s_i',s_{-i})\quad\text{for every } s_i'\in S_i$$
\bsk
Equivalent definition of Nash equilibrium:
\vfill
\centering
\large
$s$ is a Nash equilibrium \ $\iff$ \ every $s_i$ is a best reply to $s_{-i}$
\vfill
\normalsize
\raggedright
This is the version worth memorizing. It also tells you how to \textit{find} equilibria.
\vspace{-0.4in}
\end{frame}

\begin{frame}{The Underlining Trick}
\begin{enumerate}
    \item In each \textbf{column}, underline Player 1's largest payoff (all ties)
    \item In each \textbf{row}, underline Player 2's largest payoff (all ties)
    \item Any cell with \textbf{both} numbers underlined is a Nash equilibrium
\end{enumerate}
\bsk
\centering
\begin{matrixgame}[top={E, F, G, H}, left={A, B, C, D}, player names={Player 1, Player 2}]
  \payoffmatrix{4,0 & 3,2 & 2,3 & 4,1 \\ 4,2 & 2,1 & 1,2 & 0,2 \\ 3,6 & 5,5 & 3,1 & 5,0 \\ 2,3 & 3,2 & 1,2 & 3,3}
\end{matrixgame}
\bsk
\raggedright
Your turn. Columns first, then rows.
\vspace{-0.2in}
\end{frame}

\begin{frame}{The Underlining Trick}
\centering
\begin{matrixgame}[top={E, F, G, H}, left={A, B, C, D}, player names={Player 1, Player 2}, br]
  \payoffmatrix{4,0 & 3,2 & 2,3 & 4,1 \\ 4,2 & 2,1 & 1,2 & 0,2 \\ 3,6 & 5,5 & 3,1 & 5,0 \\ 2,3 & 3,2 & 1,2 & 3,3}
\end{matrixgame}
\bsk
\raggedright
Only $(B,E)$ has both payoffs underlined: the unique Nash equilibrium.
\bsk
Watch the near-misses:
\begin{itemize}
    \item $(A,E)$: Player 1 is happy, Player 2 would rather move to $G$
    \item $(C,F)$: Player 1 is happy, Player 2 would rather move to $E$
    \item $(D,H)$: Player 2 is happy, Player 1 would rather move to $C$
\end{itemize}
\vspace{-0.2in}
\end{frame}

\begin{frame}{More Than Two Players}
A profile $s^*\in S$ is a Nash equilibrium if, for every player $i$,
$$\pi_i(s^*)\ \ge\ \pi_i(s_1^*,\ldots,s_{i-1}^*,\ s_i,\ s_{i+1}^*,\ldots,s_n^*)
\quad\text{for all } s_i\in S_i$$
\bsk
$n$ inequalities, one per player. Nothing new, just more of them.
\vfill
\vspace{0.8in}
\end{frame}

\begin{frame}{A Three-Player Example}
Let's try a 3-player game ($2\times 2\times 2$)\bsk
\centering
\begin{matrixgame}[top={L, R}, left={T, B}, bottom={Player 3: $E$, Player 3: $F$}, h=0.7,
                   player names={Player 1, Player 2, Player 3}, br=<2->]
  \payoffmatrix{4,1,4 & 0,0,1 \\ 3,1,3 & 3,3,2}
  \payoffmatrix[pos={5.4,0}]{0,0,1 & 2,2,3 \\ 1,4,2 & 1,0,3}
\end{matrixgame}
\vfill
\raggedright
Two equilibria in here. Find them using the underlining technique.
\onslide<3>{
\bsk
$(T,L,E)$ and $(T,R,F)$: the only cells with all three payoffs underlined.
\par\smallskip
{\footnotesize Check $(B,R,E)$ and $(B,L,F)$: two underlines each. Players 1 and 2 are content; only Player 3 wants to move.}
}
\vspace{-0.3in}
\end{frame}
\begin{frame}{Nash Survives IDSDS}
All three procedures return a \textbf{set of strategy profiles}, not one prediction:
\begin{itemize}
    \item $NE(G)$: the Nash equilibrium profiles for game $G$
    \item $IDSDS(G)$: the profiles that survive IDSDS in game $G$
    \item $IDWDS(G)$: the profiles that survive IDWDS in game $G$
\end{itemize}
Each is a subset of $S$, and each can be empty, a single profile, or many.
\vfill
\textbf{Theorem.} \ For every ordinal strategic-form game $G$,
$$NE(G)\ \subseteq\ IDSDS(G)$$
\begin{itemize}
    \item A Nash equilibrium never uses a strategy that IDSDS deletes
    \item So IDSDS never throws away a Nash equilibrium: it's safe
    \item But the same is \textbf{not} true of IDWDS: it is possible that $NE(G)\cap IDWDS(G)=\emptyset$
\end{itemize}
\vspace{-0.3in}
\end{frame}

\begin{frame}{Strict Nash Equilibrium}
$s^*$ is a \textbf{strict} Nash equilibrium if every player is \textit{strictly} worse off
deviating, rather than merely no better off.
\bsk
\centering
\begin{tabular}{ll}
Nash: & $\pi_i(s^*)\ \ge\ \pi_i(s_i,s_{-i}^*)$ \ \ for every $s_i\in S_i$ \\[0.18in]
Strict: & $\pi_i(s^*)\ >\ \pi_i(s_i,s_{-i}^*)$ \ \ for every $s_i\ne s_i^*$ \\
\end{tabular}
\bsk
\raggedright
Two changes: $\ge$ becomes $>$, and $s_i^*$ itself is excluded. It has to be, since
$\pi_i(s^*)>\pi_i(s^*)$ is false.
\vfill
\textbf{Theorem.} \ If $IDSDS(G)=\{s^*\}$ then $s^*$ is a strict Nash equilibrium.
\bsk
If $IDWDS(G)=\{s^*\}$ then $s^*$ is a Nash equilibrium, but not necessarily a strict one.
\vspace{-0.3in}
\end{frame}

\begin{frame}{A Bigger Example}
Suppose OSU donor Ratmir Timashev makes the following offer to 50 students.
Each secretly writes down a request: a multiple of \$10, up to \$100.
\bsk
\begin{itemize}
    \item If \textbf{5 or fewer} students ask for \$100, everyone gets what they asked for
    \item Otherwise, everyone gets nothing
\end{itemize}
\bsk
All students are selfish and greedy. What are the Nash equilibria?
\vfill
\vspace{1.1in}
\end{frame}

\begin{frame}{A Bigger Example}
There are two families of equilibria.
\bsk
\textbf{(a) Seven or more students request \$100.}\\
Everyone gets nothing. Can anyone gain by switching alone?\\
Is this a \textbf{strict} Nash equilibrium?\vspace{0.4in}
\bsk
\textbf{(b) Exactly five request \$100, and all 45 others request \$90.}\\
Everyone is paid. Can anyone gain by switching alone?\\
Is this a \textbf{strict} Nash equilibrium?\vspace{0.4in}
\bsk
{\footnotesize Convince yourself that \textit{exactly six} requesting \$100 is not an equilibrium.}
\end{frame}

\begin{frame}{When There Is None}
Matching Pennies. Both match $\Rightarrow$ Player 1 wins; otherwise Player 2 wins.
\bsk
\centering
\begin{matrixgame}[top={H, T}, left={H, T}, player names={Player 1, Player 2}, br]
  \payoffmatrix{1,0 & 0,1 \\ 0,1 & 1,0}
\end{matrixgame}
\bsk
\raggedright
No cell has both payoffs underlined.
\vfill
With only ordinal payoffs, a game need not have \textbf{any} Nash equilibrium.
\bsk
Fixing this requires letting players randomize, and randomizing requires
cardinal payoffs. That's Chapter 5.
\vspace{-0.3in}
\end{frame}

% ==============================================================
% 2.7 INFINITE STRATEGY SETS
% ==============================================================

\begin{frame}[standout]
Infinite Strategy Sets
\end{frame}

\begin{frame}{No Table, Same Definition}
Two players. Each writes down a real number $\ge 1$, so $S_1=S_2=[1,\infty)$.
\bsk
Player 1 writes $x$, Player 2 writes $y$:
$$\pi_1(x,y)=\begin{cases} x-1 & \text{if } x<y\\ 0 & \text{if } x\ge y\end{cases}
\qquad
\pi_2(x,y)=\begin{cases} y-1 & \text{if } x>y\\ 0 & \text{if } x\le y\end{cases}$$
\bsk
You want to write the larger\ldots\ no, the \textit{smaller} number. But not too small.
\vfill
What are the Nash equilibria? \vspace{0.7in}
\end{frame}

\begin{frame}{The Answer Is $(1,1)$, and Only $(1,1)$}
\textbf{It is an equilibrium.} At $(1,1)$ both get 0. If Player 1 switches to any
$x>1$, then $x\ge y=1$, so they still get 0.
\bsk
\textbf{Nothing else is.}
\begin{itemize}
    \item If $x=y>1$: Player 1 gets 0, but any $\hat x$ with $1<\hat x<x$ gives
      $\hat x-1>0$
    \item If $x<y$: Player 1 gets $x-1$, but any $\hat x$ with $x<\hat x<y$ gives
      $\hat x-1>x-1$
    \item If $y<x$: same argument for Player 2
\end{itemize}
\vfill
\textbf{And now the strange part.} For Player 1, $x=1$ is \textbf{weakly dominated}
by every other strategy.
\bsk
So the unique Nash equilibrium has both players using a weakly dominated strategy.
\vspace{-0.3in}
\end{frame}

\begin{frame}{Cournot Competition}
{\footnotesize (This one uses calculus.)}
\bsk
Shell and Speedway sell an identical good: a gallon of regular. Station $i$
chooses quantity $q_i\ge 0$ at constant unit cost $c$. The price is a \textit{function}
of total output:
$$P(q_1,q_2)=a-b(q_1+q_2),\qquad a>c,\ b>0$$
\bsk
Shell's profit:
$$\pi_1(q_1,q_2)=P(q_1,q_2)\,q_1-cq_1=(a-c)q_1-b q_1^2-b q_1 q_2$$
\bsk
Strategy sets are $[0,\infty)$ (infinite). But a Nash equilibrium is still just a
pair $(\bar q_1,\bar q_2)$ where neither station wants to move.
\vfill
{\footnotesize Cournot wrote this down in 1838, inventing Nash equilibrium about
120 years early, for this one class of games.}
\vspace{-0.3in}
\end{frame}

\begin{frame}{Variables and Parameters}
Not every letter in a model plays the same role.
\bsk
\textbf{Parameters:} \ $a$, $b$, $c$
\begin{itemize}
    \item fixed numbers, known to everyone, chosen by nobody
    \item placeholders for a number that is already settled before anyone moves
    \item left as letters so one solution covers every version of the game
\end{itemize}
\bsk
\textbf{Choice variables:} \ $q_1$, $q_2$
\begin{itemize}
    \item placeholders for a number that someone gets to pick
    \item these are the strategies: $S_i=[0,\infty)$
\end{itemize}
\bsk
\textbf{Derived values} (functions of the variables): \ $P(q_1,q_2)$, \ $\pi_i(q_1,q_2)$
\begin{itemize}
    \item not fixed in advance, and nobody chooses them directly
    \item they are \textit{functions}: feed in the choices, get a number out
\end{itemize}
\vfill
For any letter, ask: could a player have made it come out differently? \vspace{0.3in}
\end{frame}

% ---- Calculus reminder: partials, FOC, SOC --------------------------
\begin{frame}{Maximizing, With Two Variables}
From the math review: to maximize $f(x)$, set $f'(x)=0$.
\bsk
Shell's profit depends on \textit{two} quantities:
$$\pi_1(q_1,q_2)=(a-c)q_1-bq_1^2-bq_1q_2$$
\bsk
But Shell only chooses $q_1$. Speedway's $q_2$ is not Shell's to pick.
\vfill
So treat $q_2$ as a fixed number and maximize over $q_1$ alone. \vspace{0.8in}
\end{frame}

\begin{frame}{Partial Derivatives}
A \textbf{partial derivative} is an ordinary derivative taken while every other variable
is held constant, treated as though it were just a number.
\bsk
$$\frac{\partial f}{\partial x}\ :\quad \text{differentiate in } x, \text{ treat } y
\text{ as a constant}$$
\bsk
The curly $\partial$ is the only new thing. The rules are the ones you already know.
\vfill
$f(x,y)=3x^2+5xy+y^3$
\bsk
$\dfrac{\partial f}{\partial x}=$ \underline{\hspace{2in}} \vspace{0.4in}

$\dfrac{\partial f}{\partial y}=$ \underline{\hspace{2in}} \vspace{0.3in}
\end{frame}

\begin{frame}{Why That Works}
Same $f(x,y)=3x^2+5xy+y^3$. Suppose $y$ is stuck at $2$:
$$f(x,2)=3x^2+10x+8$$
\bsk
That is a function of $x$ alone. Differentiate it the ordinary way:
\bsk
$\dfrac{d}{dx}f(x,2)=$ \underline{\hspace{2in}} \vspace{0.5in}

Now compare it to $\dfrac{\partial f}{\partial x}$ with $y=2$ substituted in.
\vfill
Holding a variable constant and plugging in a constant are the same operation.
\vspace{0.5in}
\end{frame}

\begin{frame}{Finding the Maximizer}
Two conditions, in words first:
\begin{itemize}
    \item \textbf{First order}: at the top, the slope in the $x$ direction is flat
    \item \textbf{Second order}: the curve bends downward there, so flat means a peak
      rather than a valley
\end{itemize}
\bsk
$$\underbrace{\frac{\partial f}{\partial x}=0}_{\text{slope is flat}}
\qquad\text{and}\qquad
\underbrace{\frac{\partial^2 f}{\partial x^2}<0}_{\text{bending downward}}$$
\vfill
Skip the second condition and a flat point could just as easily be a minimum.
\vspace{0.7in}
\end{frame}

\begin{frame}{Shell's Problem}
$$\pi_1(q_1,q_2)=(a-c)q_1-bq_1^2-bq_1q_2$$
\bsk
Hold $q_2$ fixed. First order:
\bsk
$\dfrac{\partial \pi_1}{\partial q_1}=$ \underline{\hspace{2.2in}} $=0$ \vspace{0.6in}

Second order:
\bsk
$\dfrac{\partial^2 \pi_1}{\partial q_1^2}=$ \underline{\hspace{1.2in}}, negative because
$b>0$, so the flat point is a maximum. \vspace{0.4in}
\end{frame}

\begin{frame}{From First-Order Condition to Best Reply}
Shell's first-order condition holds at Shell's optimum, \textit{whatever} $q_2$ turns out
to be:
$$a-c-2bq_1-bq_2=0$$
\bsk
So solve it for $q_1$: \vspace{0.9in}

\centering
$\displaystyle BR_1(q_2)=\frac{a-c-bq_2}{2b}$
\vfill
\raggedright
The second-order condition was $-2b<0$, so this really is Shell's best, not its worst.
\vspace{-0.2in}
\end{frame}

\begin{frame}{A Best Reply Is a Function}
$$BR_1(q_2)=\frac{a-c-bq_2}{2b}$$
\bsk
Feed it a quantity for Speedway, and it hands back Shell's best answer to that quantity.
One input, one output.
\bsk
\begin{itemize}
    \item If Speedway floods the market, $BR_1$ tells Shell to pull back
    \item If Speedway sits out, $BR_1$ tells Shell to act like a monopolist
\end{itemize}
\bsk
Speedway's problem is the same with the labels swapped:
$$BR_2(q_1)=\frac{a-c-bq_1}{2b}$$
\vfill
{\footnotesize $BR_i(q_{-i})$ is truncated at 0 if needed.}
\vspace{-0.3in}
\end{frame}

\begin{frame}{Nash Equilibrium Is a Fixed Point}
Recall the best-reply version of the definition: $(q_1^*,q_2^*)$ is a Nash equilibrium
exactly when each is a best reply to the other.
$$q_1^*=BR_1(q_2^*) \qquad\text{and}\qquad q_2^*=BR_2(q_1^*)$$
\bsk
Substitute the second into the first:
$$q_1^*=BR_1\big(BR_2(q_1^*)\big)$$
\bsk
Two equations in two unknowns just became \textbf{one} equation in one unknown.
\vfill
Start at $q_1^*$, ask what Speedway would do, then ask what Shell would do back. A Nash
equilibrium is a quantity that returns to \textit{itself} after that round trip.
\vspace{0.3in}
\end{frame}

\begin{frame}{Solving the Fixed Point}
$$q_1=\frac{a-c-b\,BR_2(q_1)}{2b},\qquad BR_2(q_1)=\frac{a-c-bq_1}{2b}$$
\bsk
Substitute, then solve for $q_1$: \vspace{2.4in}

\centering
$\displaystyle q_1^*=\frac{a-c}{3b}$
\end{frame}

\begin{frame}{And Then the Other One}
Feed $q_1^*$ back through Speedway's best reply:
$$q_2^*=BR_2(q_1^*)=\frac{a-c-b\left(\frac{a-c}{3b}\right)}{2b}$$
\bsk
$q_2^*=$ \underline{\hspace{1.6in}} \vspace{0.8in}

\raggedright
No surprise that it matches $q_1^*$: the two stations face identical problems, so the
equilibrium had to be symmetric.
\vfill
Worth checking directly that $BR_1(q_2^*)=q_1^*$. It does.
\vspace{-0.2in}
\end{frame}

\begin{frame}{Cournot: The Numbers}
At $\bar q_1=\bar q_2=\dfrac{a-c}{3b}$:
\bsk
\begin{itemize}
    \item Price: \hspace{0.55in} $\displaystyle P(\bar q_1,\bar q_2)=\frac{a+2c}{3}$
    \bsk
    \item Each profit: \hspace{0.1in} $\displaystyle \frac{(a-c)^2}{9b}$
\end{itemize}
\bsk
Try it: \ $a=25$, \ $b=2$, \ $c=1$.
\bsk
\centering
$\bar q_1=\bar q_2=$ \underline{\hspace{0.8in}} \hspace{0.4in}
$P=$ \underline{\hspace{0.8in}} \hspace{0.4in}
profit $=$ \underline{\hspace{0.8in}}
\vfill
\end{frame}

% ==============================================================
% WRAP
% ==============================================================

% ==============================================================
% WORKED EXAMPLE
% ==============================================================

\begin{frame}[standout]
Let's Work One Out
\end{frame}

\begin{frame}{Three Bars on High Street}
Ann, Ben, and Cal each run a bar. Before a home game each independently commits
to a \textbf{Big} promotion or a \textbf{Small} one.
\bsk
\centering
{\small
\begin{matrixgame}[top={Big, Small}, left={Big, Small}, bottom={Cal: Big, Cal: Small}, player names={Ann, Ben, Cal}]
  \payoffmatrix{2,2,2 & 0,3,1 \\ 3,0,1 & 1,1,0}
  \payoffmatrix[pos={5.4,0}]{1,1,3 & 0,0,2 \\ 0,2,0 & 2,2,3}
\end{matrixgame}}
\bsk
\raggedright
{\footnotesize Ann picks the row, Ben the column, Cal the table. Payoffs in order:
Ann, Ben, Cal.}
\vspace{-0.2in}
\end{frame}

\begin{frame}{Three Bars on High Street}
\begin{enumerate}
    \item Does anyone have a dominant strategy? \vspace{0.25in}
    \item Does IDSDS delete anything at all? \vspace{0.25in}
    \item Find every Nash equilibrium. Underline Ann's best payoff in each column,
      Ben's in each row, and Cal's whenever it beats her payoff in the same cell
      of the other table. \vspace{0.35in}
    \item Is either equilibrium Pareto superior to the other? \vspace{0.25in}
\end{enumerate}
\end{frame}

\begin{frame}{What Have We Learned, Charlie Brown?}
\begin{itemize}
    \item \textbf{Second-price auction}: truthful bidding is weakly dominant (Vickrey)
    \item \textbf{Pivotal mechanism}: truthful reporting is weakly dominant (Clarke)
    \item \textbf{IDSDS}: delete strictly dominated strategies, repeat (order-free,
      justified by common knowledge of rationality)
    \item \textbf{IDWDS}: order matters, so delete simultaneously, and the
      justification is murkier
    \item \textbf{Nash equilibrium}: nobody gains by deviating alone
    \item $NE\subseteq IDSDS$, but a game may have no Nash equilibrium at all
\end{itemize}
\vfill
Next: what happens when players can randomize.
\vspace{-0.4in}
\end{frame}

\begin{frame}{Reading \& Practice}
\begin{itemize}
    \item Bonanno, \textit{Game Theory}, \S2.3--2.7
    \item Exercises: \S2.9.3 (auctions), \S2.9.4 (pivotal mechanism),\\
      \S2.9.5 (iterated deletion), \S2.9.6 (Nash equilibrium), \S2.9.7 (infinite games)
    \item Proofs of both theorems are in \S2.8
    \item Solutions in \S2.10. Try them first
\end{itemize}
\vfill
\end{frame}

\end{document}
